% ============================================================
%  Semana 12 · Clase 23: Clínica de analítica predictiva y storytelling
%  Curso: Análisis y Visualización de Datos para la Toma de Decisiones
%
%  Formato de la sesión (3 h):
%    Hora 1. El pitch: estructura, apertura con problema y contexto,
%            evidencia, el modelo explicado sin jerga, recomendación e
%            impacto esperado
%    Hora 2. Incertidumbre y limitaciones; gestión de preguntas;
%            retroalimentación estructurada
%    Hora 3. Clínica: presentaciones cortas en grupos, con
%            retroalimentación estructurada
%
%  Ejemplo: el modelo de bicicletas de la Clase 20 (Bike Sharing Dataset,
%  UCI, CC BY 4.0). Las cifras son las de la práctica de esa clase:
%  train_test_split(test_size=0.2, random_state=42).
%
%  Compilar (dos veces) directamente desde esta carpeta:
%    xelatex semana12-clase23.tex
%    xelatex semana12-clase23.tex
% ============================================================
\documentclass[aspectratio=169,11pt]{beamer}

\usepackage{../../../plantilla/beamerthemeesumer}
\usetheme{esumer}

\usepackage[spanish,es-noshorthands]{babel}
\usepackage{booktabs}
\usepackage{graphicx}
\usepackage{array}
\newcolumntype{L}[1]{>{\raggedright\arraybackslash}p{#1}}
\newcolumntype{C}[1]{>{\centering\arraybackslash}p{#1}}
\usetikzlibrary{arrows.meta,positioning,calc}
\tikzset{every node/.append style={execute at begin node={\hyphenpenalty=10000\exhyphenpenalty=10000\relax}}}

% ---------------- DATOS DE LA CLASE ----------------
\title{Clínica de analítica predictiva y storytelling}
\subtitle{Presentar el caso, el modelo y la recomendación}
\author{Docente: Cristian A. Jimenez C., MSc}
\date{Nov 24 al 29}

\setmodulo{Módulo 4: Innovar (Analítica Predictiva y Storytelling)}
\setsemana{12}
\setclase{23}
\setfechaclase{Nov 24 al 29}
\setmodalidad{Virtual}
% -----------------------------------------------------

\newcommand{\urlkit}{https://github.com/craljimenez/An-lisis_y_Visualizacion_de_Datos_para_la_Toma_de_Decisiones/blob/main/clases/semana12-clase23-clinica-analitica-storytelling/informe-tecnico-kit.pdf}
\newcommand{\urlkitdatos}{https://github.com/craljimenez/An-lisis_y_Visualizacion_de_Datos_para_la_Toma_de_Decisiones/raw/main/clases/semana12-clase23-clinica-analitica-storytelling/datos/ventas_diarias_tiendas.csv}
\newcommand{\fuentebici}{\par\vspace{1pt}{\fontsize{5.5}{6.5}\selectfont\color{black!55}%
  Cifras del modelo de la Clase 20. Datos: Bike Sharing Dataset (Capital Bikeshare), UCI,
  CC BY 4.0. Fanaee-T y Gama (2013).\par}}

\begin{document}

{
\setbeamertemplate{footline}{}
\begin{frame}[plain,noframenumbering]
  \titlepage
\end{frame}
}

\begin{frame}[t]{Agenda de hoy (3 horas)}
  \begin{columns}[T,onlytextwidth]
    \column{0.32\textwidth}
    \begin{block}{Hora 1}
      \small El pitch: estructura, apertura, evidencia, el modelo sin
      jerga, recomendación e impacto.
    \end{block}
    \column{0.32\textwidth}
    \begin{block}{Hora 2}
      \small Incertidumbre y limitaciones, preguntas difíciles y cómo dar
      retroalimentación útil.
    \end{block}
    \column{0.32\textwidth}
    \begin{block}{Hora 3}
      \small La clínica: presentas tu caso en cinco minutos y recibes
      retroalimentación estructurada.
    \end{block}
  \end{columns}
  \vspace{0.3cm}
  \small En las Clases 21 y 22 armaste la historia y el argumento. Hoy los
  juntas con tu modelo en una presentación corta, la pruebas frente a otros
  y te llevas una lista concreta de mejoras.
\end{frame}

\begin{frame}[t]{¿Qué es una clínica?}
  \begin{center}
  \begin{tikzpicture}[
    bx/.style={draw=none, rounded corners=3pt, text=white, align=center,
               font=\footnotesize, minimum height=1.3cm, text width=2.9cm, inner sep=4pt},
    fl/.style={-{Stealth[length=2.2mm]}, line width=0.9pt, color=esumerNavy},
  ]
    \node[bx, fill=esumerNavy] (a) at (1.65,0) {\textbf{Presentas}\\5 minutos};
    \node[bx, fill=esumerOrange] (b) at (5.2,0) {\textbf{Te preguntan}\\3 minutos, como la gerencia};
    \node[bx, fill=esumerTeal] (c) at (8.75,0) {\textbf{Te retroalimentan}\\4 minutos, con la ficha};
    \node[bx, fill=esumerSky] (d) at (12.3,0) {\textbf{Ajustas}\\tu versión 2};
    \draw[fl] (a) -- (b); \draw[fl] (b) -- (c); \draw[fl] (c) -- (d);
  \end{tikzpicture}
  \end{center}
  \vspace{0.1cm}
  \begin{columns}[T,onlytextwidth]
    \column{0.48\textwidth}
    \begin{block}{Para qué sirve}
      \footnotesize Una presentación se mejora presentándola. En el papel
      todo parece claro; frente a alguien que no conoce el caso aparecen
      los huecos.
    \end{block}
    \column{0.48\textwidth}
    \begin{exampleblock}{Lo que se evalúa}
      \footnotesize La presentación, no a la persona. Equivocarse hoy es
      barato: para eso es la clínica.
    \end{exampleblock}
  \end{columns}
\end{frame}

% ============================================================
\section{Hora 1: El pitch}
% ============================================================
\begin{frame}[t]{El pitch en cinco minutos}
  \small Cinco partes, con un tiempo para cada una:
  \vspace{0.15cm}
  \begin{center}
  \begin{tikzpicture}[x=1cm, y=1cm]
    \foreach \x/\w/\c/\t/\s in {
      0/1.4/esumerSky/Apertura/{30 s},
      1.4/5.6/esumerNavy/Evidencia y modelo/{2 min},
      7.0/2.8/esumerTeal/Recomendación/{1 min},
      9.8/2.8/esumerOrange/Impacto/{1 min},
      12.6/1.4/esumerMagenta/Cierre/{30 s}} {
      \fill[\c] (\x,0) rectangle ++(\w,0.9);
      \node[text=white, font=\scriptsize\bfseries, align=center] at (\x+\w/2,0.45) {\t\\[-1pt]{\mdseries\s}};
    }
  \end{tikzpicture}
  \end{center}
  \footnotesize
  \renewcommand{\arraystretch}{1.25}
  \begin{tabular}{@{}L{2.8cm}L{10.7cm}@{}}
    \toprule
    \textbf{Parte} & \textbf{Qué debe quedar claro} \\
    \midrule
    Apertura & El problema, para quién y por qué importa ahora. Con una cifra. \\
    Evidencia y modelo & El hallazgo principal y qué tan bien predice el modelo, comparado con lo que se hace hoy. \\
    Recomendación & Qué hacer, quién, cuándo (Clase 21). \\
    Impacto & Cuánto vale, con su línea base y sus supuestos (Clase 22). \\
    Cierre & La decisión que pides y una frase que resuma todo. \\
    \bottomrule
  \end{tabular}
\end{frame}

\begin{frame}[t]{La apertura: el problema primero}
  \small Los primeros 30 segundos deciden si te escuchan. Empieza por el
  problema de quien decide, no por ti ni por los datos.
  \vspace{0.1cm}

  \footnotesize
  \renewcommand{\arraystretch}{1.35}
  \begin{tabular}{@{}L{6.2cm}L{7.3cm}@{}}
    \toprule
    \textbf{Así no} & \textbf{Así sí} \\
    \midrule
    ``Hola, hoy les voy a presentar mi proyecto de regresión lineal.'' & ``Cada noche decidimos cuántas bicicletas preparar para mañana, y hoy lo hacemos con el promedio.'' \\
    ``Primero, les cuento cómo limpiamos los datos.'' & ``Con el promedio nos equivocamos en unas 1700 bicicletas al día.'' \\
    ``Usamos 731 días y 13 columnas.'' & ``Con el pronóstico del clima, ese error baja a menos de la mitad.'' \\
    \bottomrule
  \end{tabular}
  \vspace{0.15cm}
  \begin{block}{La fórmula de la apertura}
    \footnotesize \textbf{Situación} (lo que pasa hoy) + \textbf{problema}
    (lo que cuesta, con una cifra) + \textbf{promesa} (lo que vas a
    proponer). Tres frases.
  \end{block}
  \fuentebici
\end{frame}

\begin{frame}[t]{El modelo, sin jerga}
  \small La gerencia no necesita saber cómo funciona una regresión, sino
  cuatro cosas:
  \vspace{0.1cm}

  \footnotesize
  \renewcommand{\arraystretch}{1.3}
  \begin{tabular}{@{}L{3.3cm}L{4.3cm}L{5.5cm}@{}}
    \toprule
    \textbf{La pregunta} & \textbf{En jerga} & \textbf{En lenguaje de negocio} \\
    \midrule
    ¿Qué predice? & Regresión lineal múltiple sobre \texttt{alquileres} & Cuántas bicicletas se alquilarán mañana \\
    ¿Con qué? & Seis variables explicativas & El pronóstico del clima y el tipo de día \\
    ¿Qué tan bien? & MAE 733,4; $R^2$ 0,789 & Se equivoca en unas 730 bicicletas al día; hoy, con el promedio, en unas 1700 \\
    ¿Qué aprendió? & Coeficiente de \texttt{lluvia}: $-1239{,}71$ & Un día con lluvia se alquilan unas 1240 bicicletas menos \\
    \bottomrule
  \end{tabular}
  \vspace{0.15cm}
  \begin{alertblock}{Siempre comparado}
    \footnotesize ``730 bicicletas de error'' no dice si es bueno o malo.
    ``730, frente a 1700 de hoy'' sí. La línea base (Clase 19) es lo que hace
    entendible tu modelo.
  \end{alertblock}
  \fuentebici
\end{frame}

\begin{frame}[t]{El modelo en una diapositiva}
  \small Así se vería la diapositiva de evidencia del caso de las bicicletas:
  \vspace{0.1cm}
  \begin{center}
  \begin{tikzpicture}
    \draw[esumerNavy, line width=1pt, rounded corners=3pt, fill=white] (0,0) rectangle (13.4,4.6);
    \fill[esumerNavy] (0,3.85) rectangle (13.4,4.6);
    \node[anchor=west, text=white, font=\small\bfseries] at (0.25,4.22) {Con el pronóstico del clima, el error diario baja de 1700 a 730 bicicletas};
    \foreach \y/\h/\c/\t/\v in {2.75/4.6/esumerGray!80!black/{Hoy: el promedio}/{1712}, 1.55/1.97/esumerOrange/{Con el modelo}/{733}} {
      \node[anchor=east, font=\scriptsize] at (2.6,\y) {\t};
      \fill[\c] (2.75,\y-0.3) rectangle ++(\h,0.6);
      \node[anchor=west, font=\scriptsize\bfseries] at (2.85+\h,\y) {\v};
    }
    \node[anchor=west, font=\scriptsize, text=black!60, text width=5.6cm, align=left] at (1.2,0.75) {Error promedio por día, en bicicletas (147 días de prueba)};
    \draw[esumerNavy!40] (8.9,0.4) -- (8.9,3.5);
    \node[anchor=north west, text width=3.9cm, align=left, font=\scriptsize] at (9.2,3.55)
      {\textbf{Lo que más mueve la demanda}\\[3pt]
       Lluvia: unas 1240 menos\\
       Cada grado más: unas 150 más\\
       2012 frente a 2011: unas 1960 más\\[5pt]
       {\color{black!60}Detalle del modelo: en el anexo}};
  \end{tikzpicture}
  \end{center}
  \footnotesize Un solo gráfico, el título dice la conclusión y lo técnico
  queda en el anexo. Las barras empiezan en cero (Clase 21).
  \fuentebici
\end{frame}

\begin{frame}[t]{Recomendación e impacto esperado}
  \begin{columns}[T,onlytextwidth]
    \column{0.48\textwidth}
    \begin{block}{Recomendación (Clase 21)}
      \footnotesize
      \textbf{Qué:} planear cada día con la predicción del modelo, no con
      el promedio.\\[2pt]
      \textbf{Quién:} el equipo de operaciones.\\[2pt]
      \textbf{Cuándo:} un piloto de un mes.\\[2pt]
      \textbf{Qué pedimos hoy:} aprobar el piloto.
    \end{block}
    \column{0.48\textwidth}
    \begin{exampleblock}{Impacto esperado (Clase 22)}
      \footnotesize
      \textbf{Línea base:} planear con el promedio, con un error de unas
      1700 bicicletas al día.\\[2pt]
      \textbf{Con el modelo:} unas 730 al día en los días de prueba.\\[2pt]
      \textbf{Cómo lo sabremos:} comparar cada día del piloto lo
      planeado con lo alquilado.
    \end{exampleblock}
  \end{columns}
  \vspace{0.15cm}
  \begin{alertblock}{Cuando no hay pesos}
    \footnotesize No siempre puedes poner el impacto en dinero: aquí no
    sabemos cuánto cuesta una bicicleta de más o de menos. Dilo, y muestra
    el impacto en la unidad que tienes (bicicletas de error), con el dato
    que faltaría para pasarlo a pesos.
  \end{alertblock}
  \fuentebici
\end{frame}

\begin{frame}[t]{Un pitch completo, en cinco minutos}
  \footnotesize
  \renewcommand{\arraystretch}{1.3}
  \begin{tabular}{@{}L{2.3cm}L{11.2cm}@{}}
    \toprule
    \textbf{Parte} & \textbf{Lo que dirías} \\
    \midrule
    Apertura & ``Cada noche decidimos cuántas bicicletas preparar. Hoy usamos el promedio y nos equivocamos en unas 1700 al día. Les proponemos una forma de bajar ese error a menos de la mitad.'' \\
    Evidencia y modelo & ``Con dos años de datos, el modelo aprendió cómo cambia la demanda con el clima y el tipo de día. Probado con días que no vio, se equivoca en unas 730 bicicletas. Lo que más pesa: la lluvia resta unas 1240 y cada grado suma unas 150.'' \\
    Recomendación & ``Proponemos un piloto de un mes: operaciones planea cada día con la predicción.'' \\
    Impacto & ``Esperamos pasar de unas 1700 a unas 730 bicicletas de error por día. Para traducirlo a pesos necesitamos saber cuánto cuesta cada bicicleta de más o de menos.'' \\
    Cierre & ``Pedimos aprobar hoy el piloto. Sabremos si funciona comparando, cada día, lo planeado con lo alquilado.'' \\
    \bottomrule
  \end{tabular}
  \fuentebici
\end{frame}

% ============================================================
\section{Hora 2: Incertidumbre, limitaciones y preguntas}
% ============================================================
\begin{frame}[t]{Comunicar la incertidumbre}
  \small Una predicción sin su error parece más segura de lo que es.
  \vspace{0.1cm}
  \begin{columns}[T,onlytextwidth]
    \column{0.5\textwidth}
    \footnotesize
    \renewcommand{\arraystretch}{1.35}
    \begin{tabular}{@{}L{2.4cm}L{4.1cm}@{}}
      \toprule
      \textbf{Así no} & \textbf{Así sí} \\
      \midrule
      ``Mañana se alquilarán 6371 bicicletas.'' & ``Esperamos unas 6400; normalmente nos equivocamos en unas 730, así que planeamos entre 5600 y 7100.'' \\
      ``El modelo tiene un $R^2$ de 0,789.'' & ``El modelo explica la mayor parte de la variación, pero no toda.'' \\
      \bottomrule
    \end{tabular}
    \column{0.46\textwidth}
    \begin{block}{Tres reglas}
      \footnotesize
      \begin{enumerate}
        \footnotesize
        \item \textbf{Redondea}: 6400, no 6371,4. Los decimales dan una
          precisión que el modelo no tiene.
        \item \textbf{Da un rango} con palabras simples: ``entre\ldots
          y\ldots''.
        \item \textbf{Di de dónde sale}: ``en los días de prueba, el error
          típico fue de 730''.
      \end{enumerate}
    \end{block}
    \scriptsize\color{black!60} El rango de $\pm$730 es una regla
    práctica, no un intervalo estadístico formal.
  \end{columns}
  \fuentebici
\end{frame}

\begin{frame}[t]{Decir las limitaciones}
  \small Si no las dices tú, alguien las encontrará y dudará de todo lo
  demás.
  \vspace{0.1cm}

  \footnotesize
  \renewcommand{\arraystretch}{1.3}
  \begin{tabular}{@{}L{2.7cm}L{10.8cm}@{}}
    \toprule
    \textbf{Limitación} & \textbf{En el caso de las bicicletas} \\
    \midrule
    De los datos & Dos años de una sola ciudad. No sabemos si vale para otra. \\
    Del modelo & Es una recta: con calor fuerte predice de más, porque la demanda baja (Clase 20). \\
    Del rango & Aprendió con días entre 2,4 y 35,3 grados. Fuera de ahí no se usa. \\
    Eventos raros & El 29 de octubre de 2012, con el huracán Sandy, predijo unas 2980 y se alquilaron 22. \\
    De causa & Muestra asociaciones, no causas (Clase 19). \\
    \bottomrule
  \end{tabular}
  \vspace{0.15cm}
  \begin{exampleblock}{Una frase que lo resume}
    \footnotesize ``El modelo sirve para planear días normales con el
    pronóstico del clima. No sirve para eventos extraordinarios ni para
    temperaturas fuera de lo que hemos visto.''
  \end{exampleblock}
  \fuentebici
\end{frame}

\begin{frame}[t]{Ni esconder ni ahogar}
  \begin{columns}[T,onlytextwidth]
    \column{0.31\textwidth}
    \begin{alertblock}{Esconder}
      \footnotesize Solo las cifras buenas, sin errores ni limitaciones.
      \vspace{3pt}

      \textbf{Riesgo:} pierdes la confianza cuando aparece el primer error
      real.
    \end{alertblock}
    \column{0.31\textwidth}
    \begin{alertblock}{Ahogar}
      \footnotesize Diez minutos de advertencias antes de decir qué
      proponer.
      \vspace{3pt}

      \textbf{Riesgo:} nadie se atreve a decidir con un modelo que tú mismo
      no defiendes.
    \end{alertblock}
    \column{0.31\textwidth}
    \begin{exampleblock}{Equilibrar}
      \footnotesize El rango del error en la evidencia, una diapositiva de
      limitaciones y el detalle en el anexo.
      \vspace{3pt}

      \textbf{Resultado:} confían en el modelo y saben cuándo no usarlo.
    \end{exampleblock}
  \end{columns}
  \vspace{0.2cm}
  \small La incertidumbre no debilita la recomendación: la hace
  \textbf{creíble}. Una decisión con un rango claro es mejor que una con
  una cifra exacta y falsa.
\end{frame}

\begin{frame}[t]{Responder preguntas}
  \begin{center}
  \begin{tikzpicture}[
    bx/.style={draw=none, rounded corners=3pt, text=white, align=center,
               font=\footnotesize, minimum height=1.3cm, text width=2.3cm, inner sep=3pt},
    fl/.style={-{Stealth[length=2.2mm]}, line width=0.9pt, color=esumerNavy},
  ]
    \node[bx, fill=esumerNavy] (a) at (1.25,0) {\textbf{Escucha}\\hasta el final};
    \node[bx, fill=esumerSky] (b) at (4.05,0) {\textbf{Reformula}\\``Si entiendo bien\ldots''};
    \node[bx, fill=esumerTeal] (c) at (6.85,0) {\textbf{Responde}\\corto, con una cifra};
    \node[bx, fill=esumerOrange] (d) at (9.65,0) {\textbf{Conecta}\\con tu mensaje};
    \node[bx, fill=esumerMagenta] (e) at (12.45,0) {\textbf{Verifica}\\``¿Responde a su pregunta?''};
    \draw[fl] (a) -- (b); \draw[fl] (b) -- (c); \draw[fl] (c) -- (d); \draw[fl] (d) -- (e);
  \end{tikzpicture}
  \end{center}
  \vspace{0.1cm}
  \begin{columns}[T,onlytextwidth]
    \column{0.48\textwidth}
    \begin{block}{Si no sabes la respuesta}
      \footnotesize ``No lo sé; lo reviso y les respondo mañana.'' Es mejor
      que inventar. Anota la pregunta y cumple.
    \end{block}
    \column{0.48\textwidth}
    \begin{block}{Si la pregunta es un ataque}
      \footnotesize No te defiendas: busca la preocupación de fondo
      (``¿le preocupa que el modelo falle en días de lluvia?'') y responde
      a eso, con datos.
    \end{block}
  \end{columns}
\end{frame}

\begin{frame}[t]{Las preguntas que siempre llegan}
  \small Prepáralas antes. La mayoría tiene su respuesta en el anexo:
  \vspace{0.1cm}

  \footnotesize
  \renewcommand{\arraystretch}{1.3}
  \begin{tabular}{@{}L{4.0cm}L{9.5cm}@{}}
    \toprule
    \textbf{Pregunta} & \textbf{Cómo prepararla} \\
    \midrule
    ¿Qué tan seguro estás? & El error típico en días que el modelo no vio, y el rango. \\
    ¿Por qué confiar en el modelo? & Comparado con lo de hoy (la línea base), y cómo se probó: entrenamiento y prueba. \\
    ¿Qué pasa si no hacemos nada? & La línea base: el error o el costo de seguir igual. \\
    ¿Cuánto cuesta y cuánto deja? & El impacto, con sus supuestos; si falta un dato, cuál. \\
    ¿Y si los datos están mal? & Qué revisaste y qué decidiste (la bitácora de la Clase 12). \\
    ¿Cuándo no funciona? & Tu diapositiva de limitaciones. \\
    \bottomrule
  \end{tabular}
\end{frame}

\begin{frame}[t]{El anexo: lo que no presentas, pero tienes listo}
  \begin{columns}[T,onlytextwidth]
    \column{0.48\textwidth}
    \small Después de la diapositiva de cierre, prepara de 3 a 6
    diapositivas de respaldo:
    \begin{itemize}
      \footnotesize
      \item Los datos: fuente, periodo, filas y la limpieza que hiciste.
      \item El modelo: variables, coeficientes y medidas de error.
      \item La comparación con la línea base, en una tabla.
      \item Los supuestos del impacto y su sensibilidad.
      \item Los casos en que más se equivoca.
    \end{itemize}
    \column{0.48\textwidth}
    \begin{exampleblock}{Para qué sirve}
      \footnotesize Cuando llegue la pregunta técnica, vas al anexo y
      respondes con una diapositiva preparada, en vez de improvisar.
    \end{exampleblock}
    \begin{block}{Y para la audiencia técnica}
      \footnotesize Si te piden el detalle, el anexo, el notebook o el flujo
      de Orange son tu respaldo (Clase 21).
    \end{block}
  \end{columns}
\end{frame}

\begin{frame}[t]{Dar retroalimentación útil}
  \begin{columns}[T,onlytextwidth]
    \column{0.48\textwidth}
    \begin{block}{Reglas}
      \footnotesize
      \begin{enumerate}
        \footnotesize
        \item Sobre la presentación, no sobre la persona.
        \item \textbf{Específica}: ``la diapositiva 3 tiene dos ideas'', no
          ``estuvo confusa''.
        \item Con lo que funcionó \textbf{y} lo que cambiarías.
        \item Una sugerencia concreta por cada problema.
        \item Pocas cosas importantes, no veinte detalles.
      \end{enumerate}
    \end{block}
    \column{0.48\textwidth}
    \begin{exampleblock}{Una fórmula}
      \footnotesize \textbf{Qué funcionó:} ``La apertura con la cifra del
      error me enganchó.''\\[3pt]
      \textbf{Qué cambiaría:} ``No entendí qué pides al final: dilo como
      una decisión.''\\[3pt]
      \textbf{Una pregunta:} ``¿Qué pasa con el modelo en días de
      calor?''
    \end{exampleblock}
    \footnotesize Y quien la recibe: escucha, anota y agradece. No se
    defiende: decide después qué cambiar.
  \end{columns}
\end{frame}

\begin{frame}[t]{La rúbrica de la clínica}
  \footnotesize
  \renewcommand{\arraystretch}{1.22}
  \begin{tabular}{@{}L{2.5cm}L{5.2cm}L{5.4cm}@{}}
    \toprule
    \textbf{Criterio} & \textbf{En proceso} & \textbf{Logrado} \\
    \midrule
    Apertura & Empieza por los datos o por el método. & Problema, cifra y promesa en 30 segundos. \\
    Mensaje & Hay que adivinar qué propone. & Una frase clara que se repite al cierre. \\
    Modelo & Jerga, sin comparación. & Qué predice, con qué, qué tan bien frente a la línea base. \\
    Incertidumbre & Cifras exactas, sin error. & Rango claro y limitaciones dichas. \\
    Recomendación & ``Hay que analizar más''. & Qué, quién, cuándo y qué se pide hoy. \\
    Impacto & Sin cifra o sin supuestos. & Estimado, con su línea base y sus supuestos. \\
    Tiempo & Pasa de cinco minutos. & Termina a tiempo. \\
    Preguntas & Se defiende o improvisa. & Reformula, responde corto, usa el anexo. \\
    \bottomrule
  \end{tabular}
\end{frame}

% ============================================================
\section{Hora 3: La clínica}
% ============================================================
\begin{frame}[t]{La dinámica de la clínica}
  \begin{columns}[T,onlytextwidth]
    \column{0.55\textwidth}
    \small En \textbf{grupos de tres} (salas separadas en la reunión
    virtual). Cada uno presenta su caso o el del kit. Cada ronda tiene tres
    papeles:
    \vspace{0.05cm}

    \footnotesize
    \renewcommand{\arraystretch}{1.25}
    \begin{tabular}{@{}L{1.9cm}L{5.3cm}@{}}
      \toprule
      \textbf{Papel} & \textbf{Qué hace} \\
      \midrule
      Presentador & Cuenta su caso en 5 minutos. \\
      Gerente & Escucha como quien decide y hace las preguntas: 3 minutos. \\
      Observador & Llena la ficha con la rúbrica y la comparte: 4 minutos. \\
      \bottomrule
    \end{tabular}
    \vspace{0.1cm}

    \small Tres rondas: todos pasan por los tres papeles. Unos 40 minutos
    en total.
    \column{0.41\textwidth}
    \begin{block}{Si son dos}
      \footnotesize El que escucha hace de gerente y de observador a la
      vez.
    \end{block}
    \begin{block}{Si trabajas solo}
      \footnotesize Grábate cinco minutos, mírate al día siguiente y llena
      la ficha como observador. Pide a alguien de otro grupo que haga de
      gerente.
    \end{block}
    \begin{exampleblock}{Después}
      \footnotesize De vuelta a la sala principal, algunos voluntarios
      presentan su versión mejorada.
    \end{exampleblock}
  \end{columns}
\end{frame}

\begin{frame}[t]{¿No tienes caso propio? El kit de Cafés Altavista}
  \begin{columns}[T,onlytextwidth]
    \column{0.5\textwidth}
    \small Si no hiciste las actividades de las Clases 20 a 22, o tu proyecto
    no tiene un modelo, usa el kit. Está en la carpeta de la clase:
    \vspace{0.05cm}
    \begin{itemize}
      \footnotesize
      \item \href{\urlkit}{\color{esumerOrange}\texttt{informe-tecnico-kit.pdf}}:
        instrucciones y el \textbf{informe técnico} que el equipo de
        analítica le entregó a la gerencia.
      \item \href{\urlkitdatos}{\color{esumerOrange}\texttt{datos/ventas\_diarias\_tiendas.csv}}:
        los datos del informe, por si quieres explorarlos (opcional).
    \end{itemize}
    \vspace{0.05cm}
    \footnotesize\textbf{El caso:} pronosticar las ventas de cada tienda,
    cada día, con los datos de Cafés Altavista de la Clase 16: 2005
    registros de tienda y día.
    \column{0.46\textwidth}
    \begin{alertblock}{Está en jerga, a propósito}
      \footnotesize El informe habla de MAE, $R^2$, coeficientes e
      interacciones, y \textbf{no dice qué hacer}. Tu trabajo es
      traducirlo: encontrar el hallazgo, proponer la decisión y contarlo en
      cinco minutos.
    \end{alertblock}
    \begin{block}{Una pista}
      \footnotesize Mira la figura de la sección 3 del informe: una tienda
      no se comporta como las demás los fines de semana.
    \end{block}
  \end{columns}
\end{frame}

\begin{frame}[t]{La ficha de retroalimentación}
  \footnotesize
  \textbf{Presentador:} \hrulefill\quad \textbf{Observador:} \hrulefill
  \vspace{0.1cm}

  \renewcommand{\arraystretch}{1.3}
  \begin{tabular}{@{}L{2.8cm}C{1.3cm}C{1.3cm}L{7.3cm}@{}}
    \toprule
    \textbf{Criterio} & \textbf{En proceso} & \textbf{Logrado} & \textbf{Comentario específico} \\
    \midrule
    Apertura & $\square$ & $\square$ & \\
    Mensaje & $\square$ & $\square$ & \\
    Modelo & $\square$ & $\square$ & \\
    Incertidumbre & $\square$ & $\square$ & \\
    Recomendación e impacto & $\square$ & $\square$ & \\
    Tiempo y preguntas & $\square$ & $\square$ & \\
    \bottomrule
  \end{tabular}
  \vspace{0.1cm}

  \textbf{Qué funcionó:} \hrulefill\par
  \textbf{Qué cambiaría:} \hrulefill\par
  \textbf{Una pregunta:} \hrulefill
\end{frame}

\begin{frame}[t]{Producto esperado}
  \begin{columns}[T,onlytextwidth]
    \column{0.55\textwidth}
    \small Al final de la sesión, cada persona entrega:
    \begin{enumerate}
      \footnotesize
      \item La \textbf{versión 1} de su presentación, la que llevó a la
        clínica.
      \item Las \textbf{fichas de retroalimentación} que recibió.
      \item Una \textbf{lista de cambios}: qué va a ajustar y por qué, a
        partir de las fichas.
      \item La \textbf{versión 2}, con los cambios hechos.
    \end{enumerate}
    \column{0.41\textwidth}
    \begin{exampleblock}{Si tu caso es el proyecto}
      \footnotesize Usa tu historia (Clase 21), tu matriz (Clase 22) y tu
      modelo (Clase 20, o el que hayas hecho con tus datos).
    \end{exampleblock}
    \begin{block}{Si no tienes un modelo propio}
      \footnotesize Usa el kit de Cafés Altavista: presentas el informe
      técnico traducido para la gerencia, con tu propia recomendación.
    \end{block}
  \end{columns}
\end{frame}

% ============================================================
\section{Cierre}
% ============================================================
\begin{frame}[t]{Cierre y próxima sesión}
  \begin{itemize}
    \small
    \item Cinco partes en cinco minutos: apertura, evidencia y modelo,
      recomendación, impacto y cierre.
    \item La apertura empieza por el problema de quien decide, con una
      cifra.
    \item El modelo se explica sin jerga y siempre frente a la línea base.
    \item Un rango honesto y las limitaciones dichas hacen creíble la
      recomendación.
    \item A las preguntas: escucha, reformula, responde corto y usa el anexo.
  \end{itemize}
  \vspace{0.15cm}
  \begin{block}{Próxima sesión: Clase 24}
    \small MasterClass presencial: una guía para recorrer todo el Módulo 4
    con un caso nuevo, de la pregunta predictiva a la recomendación.
  \end{block}
\end{frame}

{
\setbeamertemplate{footline}{}
\begin{frame}[plain,noframenumbering]
  \begin{tikzpicture}[remember picture,overlay]
    \fill[esumerNavy] (current page.south west) rectangle (current page.north east);
    \node at (current page.center) {\color{white}\Huge\bfseries ¡Gracias!};
  \end{tikzpicture}
\end{frame}
}

\end{document}
